\documentclass[12pt]{article}
\usepackage[utf8]{inputenc}
\usepackage[T1]{fontenc}
\usepackage{hyperref}
\title{Tampulma Language Resources for K-12 Curriculum Development}
\author{Compiled by Claude Sonnet 5 (claude-sonnet-5) for Andr\'as Kornai}
\date{}
\begin{document}
\maketitle
\section{Introduction}
Tampulma (ISO 639-3 tpm, Glottocode tamp1252) is a Gur (Mabia) language of the Savannah Region of Ghana. It is also called Tampulma or Tamprusi after the ethnic group. This report is based on a single round of web searches.
\section{Primary Educational Resources}
Citations seen in search results (no online copies were found): a 1999 GILLBT reader, \emph{Tampulma mura ton 3}, and a Ghana Institute of Linguistics English--Tampulma course, \emph{Tampulma awa siliminga sema} (1977). They are cited in the Glottolog reference records below, so their availability is unconfirmed.
\section{Linguistic and Reference Materials}
Tampulma--English dictionary on Webonary (about 1,900 entries, compiled by the GILLBT Tampulma Language Project): \url{https://www.webonary.org/tampulma/}

Gray and Gray (1969), \emph{Tampulma clauses}, Ghana Institute of Linguistics (bibliographic record only): \url{https://glottolog.org/resource/reference/id/452308}

Bergman (1964), unpublished grammatical sketch (bibliographic record only): \url{https://glottolog.org/resource/reference/id/12027}

Bibliographic record for the 1977 course: \url{https://glottolog.org/resource/reference/id/49658}

Wikipedia article (titled Tamprusi language): \url{https://en.wikipedia.org/wiki/Tamprusi_language}
\section{Religious and Community Resources}
None were found in the search results.
\section{Digital and Audio Resources}
The Webonary dictionary is the only digital resource found. No audio was confirmed.
\section{Curriculum Development Context}
A usable dictionary and older GILLBT literacy materials exist, but no school programme was found. Any curriculum would need to build on the dictionary and on locating or republishing the GILLBT readers.
\end{document}